%============================================================
% CHƯƠNG 5: CÁC GIẢI PHÁP VÀ ĐÓNG GÓP NỔI BẬT
%============================================================
\chapter{CÁC GIẢI PHÁP VÀ ĐÓNG GÓP NỔI BẬT}
\label{chap:contribution}

Chương này tập trung làm nổi bật ba giải pháp kỹ thuật đặc sắc và có
giá trị thực tiễn cao nhất trong đồ án. Với mỗi giải pháp, cấu trúc
trình bày gồm ba phần: đặt vấn đề (bài toán kỹ thuật cần giải quyết),
giải pháp (cách tiếp cận và triển khai cụ thể), và kết quả đạt được.
Ba đóng góp nổi bật bao gồm: (1) pipeline dự báo chuỗi thời gian LSTM
phục vụ qua FastAPI microservice, (2) nội suy không gian Kriging/IDW
để tạo heatmap phân bố môi trường thời gian thực, và (3) kiến trúc
truyền thông MQTT kết hợp Socket.io đảm bảo độ trễ thấp cho toàn hệ
thống.

%------------------------------------------------------------
\section{Pipeline dự báo chuỗi thời gian LSTM qua FastAPI
microservice}
\label{sec:contrib:lstm}
%------------------------------------------------------------

\subsection{Đặt vấn đề}
\label{subsec:contrib:lstm:problem}

Trong các hệ thống quan trắc môi trường truyền thống, dữ liệu cảm biến
chỉ được thu thập và hiển thị theo thời gian thực mà không có khả năng
dự báo xu hướng tương lai. Điều này đặt người vận hành vào thế bị động
- chỉ phản ứng khi các chỉ số đã vượt ngưỡng nguy hiểm thay vì chủ
động can thiệp trước.

Bài toán đặt ra là tích hợp khả năng dự báo chuỗi thời gian vào hệ
thống với các ràng buộc kỹ thuật cụ thể:

\begin{itemize}[leftmargin=2cm]
  \item \textbf{Không block luồng chính:} Backend Node.js đang xử lý
    đồng thời MQTT, Socket.io và REST API. Việc chạy inference ML nặng
    trong cùng tiến trình sẽ block event loop, gây gián đoạn toàn bộ
    hệ thống.
  \item \textbf{Tích hợp PyTorch:} Node.js không có binding native
    tốt cho PyTorch. Chạy Python subprocess trong Node.js là giải pháp
    tạm thời, không ổn định và khó mở rộng.
  \item \textbf{Chuẩn hóa dữ liệu đầu vào:} Mô hình LSTM yêu cầu dữ
    liệu đầu vào được chuẩn hóa (normalize) bằng cùng scaler đã dùng
    khi huấn luyện. Việc tái tạo scaler mỗi lần inference là tốn kém
    và dễ sai.
  \item \textbf{Khả năng mở rộng độc lập:} Khi cần thay thế mô hình
    ML (ví dụ từ LSTM sang Transformer) không nên ảnh hưởng đến
    backend Node.js.
\end{itemize}

\subsection{Giải pháp}
\label{subsec:contrib:lstm:solution}

Hệ thống giải quyết các ràng buộc trên bằng cách tách ML inference
thành \textbf{hai FastAPI microservice độc lập}:

\begin{itemize}[leftmargin=2cm]
  \item \textbf{Time-Series Forecast service} (cổng 8000, file
    \texttt{forecast/time\_forecast/main.py}): Nhận request từ
    Node.js, query PostgreSQL, chạy LSTM inference và trả về mảng dự
    báo 6/10/15 phút hoặc 1/3/6/12/24 giờ.
  \item \textbf{Spatial Forecast service} (cổng 8001, file
    \texttt{forecast/spatial\_forecast.py}): Chạy LSTM cho từng
    node rồi cấp dữ liệu cho Kriging/IDW nội suy lưới 40$\times$50,
    trả về dữ liệu heatmap cho Leaflet.
\end{itemize}

Kiến trúc này áp dụng nguyên tắc \textit{separation of concerns}:
Node.js chịu trách nhiệm về luồng dữ liệu và nghiệp vụ, Python
service chịu trách nhiệm về tính toán ML.

\textbf{Quy trình huấn luyện mô hình:}

Mã giả của quy trình huấn luyện LSTM được trình bày như sau:

\begin{lstlisting}[language=Python, caption={LSTM training pseudocode}, label={lst:lstm_train}]
# 1. Load data from PostgreSQL
data = query("SELECT temperature, humidity, aqi
              FROM readings ORDER BY timestamp")

# 2. Min-Max normalization, save scaler for inference reuse
scaler = MinMaxScaler()
data_scaled = scaler.fit_transform(data)
joblib.dump(scaler, 'scaler.pkl')

# 3. Create (X, y) pairs using sliding window
# X shape: (N, window=30, features=3)
# y shape: (N, horizon=15, features=3)
X, y = create_sequences(data_scaled,
                         window=30, horizon=15)

# 4. Define LSTM architecture (PyTorch)
class LSTMModel(nn.Module):
    def __init__(self):
        super().__init__()
        self.lstm = nn.LSTM(input_size=3,
                            hidden_size=64,
                            num_layers=2,
                            batch_first=True,
                            dropout=0.2)
        self.fc = nn.Linear(64, 15 * 3)

    def forward(self, x):
        out, _ = self.lstm(x)
        # Take hidden state at last step
        return self.fc(out[:, -1, :]).view(-1, 15, 3)

# 5. Train with Adam optimizer, MSE loss
model = LSTMModel()
optimizer = torch.optim.Adam(model.parameters(), lr=1e-3)
criterion = nn.MSELoss()

for epoch in range(num_epochs):
    pred = model(X_train)
    loss = criterion(pred, y_train)
    optimizer.zero_grad()
    loss.backward()
    optimizer.step()

# 6. Save model weights
torch.save(model.state_dict(), 'lstm_model.pth')
\end{lstlisting}

\textbf{Quy trình inference tại FastAPI service:}

\begin{lstlisting}[language=Python, caption={Inference pseudocode at FastAPI service}, label={lst:lstm_infer}]
# Startup: load model and scaler once
model = LSTMModel()
model.load_state_dict(torch.load('lstm_model.pth'))
model.eval()
scaler = joblib.load('scaler.pkl')

@app.get("/forecast")
async def forecast(device_id: int, db=Depends(get_db)):
    # 1. Fetch 30 latest records from PostgreSQL
    rows = await db.fetch(
        "SELECT temperature, humidity, aqi FROM readings
         WHERE device_id=$1
         ORDER BY timestamp DESC LIMIT 30", device_id)

    # 2. Normalize using saved scaler
    data = scaler.transform(rows)   # shape: (30, 3)

    # 3. Create input tensor
    X = torch.tensor(data).unsqueeze(0).float()
                                    # shape: (1, 30, 3)
    # 4. Inference (no_grad to save memory)
    with torch.no_grad():
        pred = model(X)             # shape: (1, 15, 3)

    # 5. Inverse transform to real values
    result = scaler.inverse_transform(
        pred.squeeze(0).numpy())    # shape: (15, 3)

    return {"forecast": result.tolist()}
\end{lstlisting}

Hình~\ref{fig:lstm_pipeline} minh họa toàn bộ pipeline từ dữ liệu
cảm biến đến kết quả dự báo hiển thị trên dashboard.

\begin{figure}[H]
  \centering
  \includegraphics[width=0.9\textwidth]{lstm_pipeline}
  \caption{Pipeline dự báo chuỗi thời gian LSTM qua FastAPI microservice}
  \label{fig:lstm_pipeline}
\end{figure}

Đồ án triển khai \textbf{hai mô hình LSTM song song}:
\texttt{lstm\_model.pth} dự báo ngắn hạn 15 phút và
\texttt{lstm\_model\_long.pth} dự báo dài hạn hơn, mỗi mô hình có
scaler riêng (\texttt{scaler.pkl} và \texttt{scaler\_long.pkl}). Người
dùng có thể chọn chế độ dự báo phù hợp từ dashboard.

\subsection{Kết quả đạt được}
\label{subsec:contrib:lstm:result}

Giải pháp FastAPI microservice mang lại các kết quả cụ thể:

\begin{itemize}[leftmargin=2cm]
  \item \textbf{Tách biệt hoàn toàn:} Backend Node.js không bị block
    trong quá trình inference. Khi FastAPI service gặp sự cố, backend
    trả về lỗi graceful thay vì crash toàn hệ thống.
  \item \textbf{Thời gian inference:} Mỗi lần dự báo 15 phút cho một
    node hoàn thành trong dưới 200ms trên môi trường CPU (không GPU),
    đáp ứng yêu cầu realtime của hệ thống.
  \item \textbf{Mô hình hội tụ tốt:} Loss curve
    (Hình~\ref{fig:lstm_loss_curve}) cho thấy training loss và
    validation loss giảm đều trong quá trình huấn luyện. Kết quả hiện
    tại được sử dụng làm mốc đánh giá ban đầu và có thể tiếp tục cải
    thiện khi mở rộng tập dữ liệu hoặc tinh chỉnh siêu tham số.
  \item \textbf{Khả năng mở rộng:} Việc thay thế kiến trúc mô hình
    chỉ cần sửa file Python trong FastAPI service mà không cần thay
    đổi bất kỳ dòng code nào trong Node.js.
\end{itemize}

\subsection{Đánh giá độ chính xác mô hình}
\label{subsec:contrib:lstm:evaluation}

Độ chính xác của mô hình được đánh giá bằng sai số tuyệt đối trung
bình (Mean Absolute Error - MAE) tại ba mốc dự báo 6 phút, 10 phút và
15 phút. Các chỉ số trong Bảng~\ref{tab:lstm_mae_by_horizon} là kết
quả thực nghiệm hiện tại, được dùng làm baseline ban đầu cho hệ thống.
Trong các lần huấn luyện tiếp theo, kết quả này có thể tiếp tục được
cải thiện bằng cách bổ sung dữ liệu huấn luyện, điều chỉnh kích thước
window, số tầng LSTM hoặc các siêu tham số tối ưu hóa.

\begin{figure}[H]
  \centering
  \includegraphics[width=0.82\textwidth]{loss_curve}
  \caption{Đường cong loss trong quá trình huấn luyện mô hình LSTM}
  \label{fig:lstm_loss_curve}
\end{figure}

\begin{table}[H]
  \centering
  \caption{Sai số MAE của mô hình LSTM theo từng mốc dự báo}
  \label{tab:lstm_mae_by_horizon}
  \begin{tabular}{lccc}
    \toprule
    \textbf{Mốc dự báo} & \textbf{MAE nhiệt độ ($^\circ$C)}
      & \textbf{MAE độ ẩm (\%)} & \textbf{MAE AQI} \\
    \midrule
    6 phút  & 0.216 & 0.910 & 4.854 \\
    10 phút & 0.244 & 0.880 & 4.921 \\
    15 phút & 0.278 & 0.921 & 5.012 \\
    \bottomrule
  \end{tabular}
\end{table}

Kết quả cho thấy sai số nhiệt độ duy trì dưới 0.3$^\circ$C ở cả ba
mốc dự báo, trong khi sai số độ ẩm xấp xỉ 1\%. Sai số AQI dao động
trong khoảng 4.854 đến 5.012, phù hợp cho mục tiêu cảnh báo xu hướng
và hỗ trợ người dùng ra quyết định sớm. Khi horizon tăng từ 6 phút lên
15 phút, MAE nhiệt độ và AQI tăng nhẹ, phản ánh đặc trưng thường gặp
của bài toán dự báo chuỗi thời gian: mốc dự báo càng xa thì độ bất
định càng lớn.

%------------------------------------------------------------
\section{Nội suy không gian Kriging/IDW để tạo heatmap phân bố
môi trường}
\label{sec:contrib:spatial}
%------------------------------------------------------------

\subsection{Đặt vấn đề}
\label{subsec:contrib:spatial:problem}

Với $n$ node cảm biến đặt rải rác trong không gian phòng, hệ thống chỉ
có $n$ điểm dữ liệu rời rạc tại $n$ vị trí cố định. Người vận hành
muốn biết điều kiện môi trường tại \textbf{bất kỳ điểm nào} trong
phòng, không chỉ tại vị trí đặt cảm biến. Các thách thức cụ thể:

\begin{itemize}[leftmargin=2cm]
  \item \textbf{Số lượng node hạn chế:} Trong môi trường thực tế,
    không thể đặt cảm biến dày đặc khắp nơi do chi phí và hạ tầng.
    Với ít node, nội suy biên (vùng góc phòng xa tất cả node) thường
    cho kết quả kém chính xác.
  \item \textbf{Phân bố không đều:} Các node thường tập trung ở một
    số vị trí thuận tiện, không phân bố đều trong không gian, gây ra
    hiệu ứng clustering làm méo kết quả nội suy.
  \item \textbf{Hiệu năng realtime:} Tính toán nội suy trên lưới
    40$\times$50 = 2000 điểm phải đủ nhanh để dashboard cập nhật
    trong vài giây.
\end{itemize}

\subsection{Giải pháp}
\label{subsec:contrib:spatial:solution}

Hệ thống triển khai chiến lược nội suy hai tầng kết hợp
\textbf{IDW để sinh node ảo} và \textbf{Kriging để nội suy chính}.

\textbf{Tầng 1 - Sinh node ảo bằng IDW:}

Để cải thiện độ bao phủ tại các góc phòng, trước khi chạy Kriging,
hệ thống dùng IDW để sinh thêm các \textbf{node ảo (virtual nodes)}
tại các góc và cạnh phòng. Giá trị của node ảo được ước tính từ các
node thực theo công thức IDW:

\begin{equation}
  \hat{z}_{corner} = \frac{\displaystyle\sum_{i=1}^{n}
    \frac{z_i}{d_i^2}}
    {\displaystyle\sum_{i=1}^{n} \frac{1}{d_i^2}}
\end{equation}

Trong đó $d_i$ là khoảng cách từ góc phòng đến node thực $i$, và
$z_i$ là giá trị nhiệt độ, độ ẩm hoặc AQI của node $i$. Các node ảo này
được thêm vào tập điểm đầu vào cho Kriging, giúp thuật toán có thêm
thông tin ở vùng biên, tránh hiện tượng extrapolation không kiểm soát.

\textbf{Tầng 2 - Nội suy Kriging trên toàn bộ lưới:}

Mã giả của quy trình nội suy không gian được trình bày dưới đây:

\begin{lstlisting}[language=Python, caption={Spatial interpolation pseudocode (Kriging/IDW)}, label={lst:spatial}]
from pykrige.ok import OrdinaryKriging
from scipy.interpolate import griddata
import numpy as np

def spatial_forecast(nodes: list[dict], room: dict):
    """
    nodes: [{"lat": ..., "lng": ...,
              "temperature": ..., "humidity": ..., "aqi": ...}, ...]
    room:  {"center_lat", "center_lng",
             "width_m", "length_m"}
    """
    # 1. Create virtual nodes at 4 room corners via IDW
    corners = get_room_corners(room)
    virtual_nodes = [idw_estimate(c, nodes)
                     for c in corners]
    all_nodes = nodes + virtual_nodes

    # 2. Extract coordinates and values
    lats = [n["lat"] for n in all_nodes]
    lngs = [n["lng"] for n in all_nodes]
    temps = [n["temperature"] for n in all_nodes]
    hums  = [n["humidity"]    for n in all_nodes]

    # 3. Create 40x50 target grid
    grid_lat = np.linspace(min(lats), max(lats), 40)
    grid_lng = np.linspace(min(lngs), max(lngs), 50)
    grid_lat2d, grid_lng2d = np.meshgrid(
        grid_lat, grid_lng, indexing='ij')

    # 4. Kriging interpolation for temperature
    ok_temp = OrdinaryKriging(
        lngs, lats, temps,
        variogram_model='gaussian',
        verbose=False, enable_plotting=False)
    z_temp, _ = ok_temp.execute(
        'grid', grid_lng, grid_lat)

    # 5. Kriging interpolation for humidity
    ok_hum = OrdinaryKriging(
        lngs, lats, hums,
        variogram_model='gaussian',
        verbose=False, enable_plotting=False)
    z_hum, _ = ok_hum.execute(
        'grid', grid_lng, grid_lat)

    return {
        "grid_lat": grid_lat.tolist(),
        "grid_lng": grid_lng.tolist(),
        "temperature_grid": z_temp.tolist(),
        "humidity_grid":    z_hum.tolist()
    }
\end{lstlisting}

Kết quả lưới 40$\times$50 điểm được gửi về frontend dưới dạng JSON.
Dashboard React dùng \textbf{react-leaflet} để render lưới này thành
lớp heatmap màu sắc (gradient từ xanh lam đến đỏ) phủ lên bản đồ sàn
phòng cho từng chỉ số nhiệt độ, độ ẩm hoặc AQI, như minh họa trong Hình~\ref{fig:heatmap_demo}.

\begin{figure}[H]
  \centering
  \includegraphics[width=0.9\textwidth]{system_spatial_forecast}
  \caption{Heatmap phân bố môi trường không gian được tạo từ nội suy Kriging, hiển thị trên bản đồ Leaflet}
  \label{fig:heatmap_demo}
\end{figure}

\subsection{Kết quả đạt được}
\label{subsec:contrib:spatial:result}

\begin{itemize}[leftmargin=2cm]
  \item \textbf{Bao phủ toàn không gian:} Từ $n$ điểm đo rời rạc, hệ
    thống tái tạo bản đồ phân bố liên tục trên toàn bộ diện tích
    phòng với độ phân giải lưới 40$\times$50 điểm.
  \item \textbf{Node ảo cải thiện biên:} Kỹ thuật sinh node ảo tại
    góc phòng bằng IDW giúp Kriging cho kết quả mượt mà hơn ở vùng
    biên, tránh hiệu ứng extrapolation bất thường khi không có node
    thực tại vùng đó.
  \item \textbf{Hiệu năng đủ dùng:} Thời gian tính toán Kriging cho
    lưới 40$\times$50 trên CPU thông thường dưới 2 giây, đáp ứng yêu
    cầu cập nhật heatmap theo yêu cầu của người dùng.
  \item \textbf{Kết hợp dự báo:} Heatmap có thể được tạo không chỉ
    từ dữ liệu thực đo hiện tại mà còn từ kết quả dự báo LSTM, cho
    phép người dùng xem \textbf{heatmap dự báo} - phân bố môi
    trường ước tính trong 15 phút tới.
\end{itemize}

%------------------------------------------------------------
\section{Kiến trúc truyền thông MQTT kết hợp Socket.io đảm bảo
độ trễ thấp}
\label{sec:contrib:comm}
%------------------------------------------------------------

\subsection{Đặt vấn đề}
\label{subsec:contrib:comm:problem}

Hệ thống quan trắc realtime đặt ra yêu cầu nghiêm ngặt về độ trễ
end-to-end: từ khi cảm biến đọc giá trị đến khi biểu đồ trên dashboard
được cập nhật phải dưới 2 giây. Cách tiếp cận đơn giản nhất là
\textbf{HTTP polling} - frontend định kỳ gọi API để lấy dữ liệu mới
- nhưng có nhiều nhược điểm:

\begin{itemize}[leftmargin=2cm]
  \item \textbf{Độ trễ không xác định:} Nếu polling mỗi 5 giây, độ
    trễ trung bình là 2.5 giây, tệ nhất là 5 giây.
  \item \textbf{Lãng phí tài nguyên:} Hầu hết các request polling
    trả về ``không có dữ liệu mới'', tốn băng thông và CPU server
    vô ích.
  \item \textbf{Không scalable:} Với 100 client đang mở dashboard,
    server nhận 20 request/giây chỉ để kiểm tra dữ liệu mới.
\end{itemize}

Ngoài ra, giao tiếp giữa node IoT và server cũng cần được thiết kế
cẩn thận: HTTP POST từ ESP32 mỗi 5 giây tạo ra overhead lớn (TCP
handshake, HTTP headers ~500 bytes) cho thiết bị có tài nguyên hạn chế.

\subsection{Giải pháp}
\label{subsec:contrib:comm:solution}

Hệ thống áp dụng chiến lược \textbf{hai giao thức bổ trợ nhau}: MQTT
cho luồng dữ liệu từ node IoT lên server, và Socket.io cho luồng dữ
liệu từ server xuống dashboard.

\textbf{Tầng 1 - MQTT (Node IoT $\to$ Server):}

MQTT được chọn vì overhead tối thiểu: header chỉ 2 byte cố định, kết
nối TCP được duy trì liên tục (persistent connection) nên không tốn
chi phí handshake mỗi lần gửi. Mỗi 5 giây, ESP32 publish một message
nhỏ gọn:

\begin{lstlisting}[language=C, caption={MQTT payload from ESP32 node},
label={lst:mqtt_payload}]
// JSON payload sent via MQTT: ~60 bytes
StaticJsonDocument<128> doc;
doc["temperature"] = temperature;
doc["humidity"]    = humidity;
doc["aqi"]         = aqi; // bat buoc cho ML pipeline

char payload[128];
serializeJson(doc, payload);

// Publish to topic device/{api_key}/data
client.publish(topic.c_str(), payload.c_str());
\end{lstlisting}

\textbf{Tầng 2 - Aedes Broker (Server-side MQTT):}

Thay vì triển khai Mosquitto broker riêng biệt, hệ thống tích hợp
\textbf{Aedes} - MQTT broker thuần Node.js - trực tiếp trong tiến
trình backend. Điều này cho phép xử lý message MQTT trong cùng event
loop với logic nghiệp vụ, tránh độ trễ mạng giữa broker và backend:

\begin{lstlisting}[language=JavaScript, caption={MQTT message handling
in Aedes broker}, label={lst:aedes}]
aedes.on('publish', async (packet, client) => {
  const topic = packet.topic;

  // Only handle sensor data topics
  if (!topic.startsWith('device/')) return;

  // 1. Extract api_key from topic
  const apiKey = topic.split('/')[1];

  // 2. Validate and get device_id from PostgreSQL
  const device = await db.query(
    'SELECT id, owner_user_id FROM devices
     WHERE api_key=$1', [apiKey]);
  if (!device.rows.length) return;  // skip if invalid key

  // 3. Parse JSON payload
  const { temperature, humidity, aqi } = JSON.parse(
    packet.payload.toString());

  // 4. Save to PostgreSQL
  await db.query(
    'INSERT INTO readings
     (device_id, temperature, humidity, aqi, timestamp)
     VALUES ($1,$2,$3,$4,NOW())',
    [device.rows[0].id, temperature, humidity, aqi]);

  // 5. Check alert thresholds (alertService)
  await alertService.check(device.rows[0], {
    temperature, humidity, aqi });

  // 6. Push realtime to dashboard via Socket.io
  io.to(`device_${device.rows[0].id}`)
    .emit('new_reading', { temperature, humidity, aqi,
                           timestamp: new Date() });
});
\end{lstlisting}

\textbf{Tầng 3 - Socket.io (Server $\to$ Dashboard):}

Khi dashboard mở, client xác thực bằng JWT và đăng ký nhận dữ liệu
của thiết bị cụ thể:

\begin{lstlisting}[language=JavaScript, caption={Client-side Socket.io in React dashboard}, label={lst:socketio_client}]
// Connect and authenticate
const socket = io(SERVER_URL);
socket.emit('authenticate', { token: jwtToken });

// Subscribe to device data
socket.emit('subscribe_device', { deviceId });

// Listen for new data and update chart
socket.on('new_reading', (data) => {
  setChartData(prev => ({
    labels: [...prev.labels, data.timestamp],
    temperature: [...prev.temperature, data.temperature],
    humidity:    [...prev.humidity,    data.humidity],
    aqi:         [...prev.aqi,         data.aqi]
  }));
});
\end{lstlisting}

Hình~\ref{fig:comm_flow} minh họa toàn bộ luồng truyền thông từ node
ESP32 đến dashboard với các mốc thời gian tương ứng.

\begin{figure}[H]
  \centering
  \includegraphics[width=0.9\textwidth]{comm_flow}
  \caption{Luồng truyền thông end-to-end của hệ thống}
  \label{fig:comm_flow}
\end{figure}

\subsection{Kết quả đạt được}
\label{subsec:contrib:comm:result}

\begin{itemize}[leftmargin=2cm]
  \item \textbf{Độ trễ end-to-end thấp:} Từ khi ESP32 publish message
    MQTT đến khi biểu đồ Chart.js trên dashboard cập nhật, độ trễ
    trung bình đo được dưới 500ms trong điều kiện mạng LAN, đáp ứng
    tốt yêu cầu realtime.
  \item \textbf{Không polling:} Socket.io push-based hoàn toàn loại bỏ
    HTTP polling, giảm số lượng request xuống backend từ $O(n_{client}
    \times f_{poll})$ về 0 khi không có dữ liệu mới.
  \item \textbf{Aedes tích hợp đơn giản hóa triển khai:} Không cần
    cài đặt và vận hành Mosquitto riêng biệt. Toàn bộ hệ thống backend
    chỉ cần một tiến trình Node.js duy nhất với lệnh \texttt{node
    server/index.js}.
  \item \textbf{Xác thực hai lớp:} Node IoT xác thực bằng
    \texttt{api\_key} qua MQTT, người dùng dashboard xác thực bằng JWT
    qua Socket.io. Hai lớp xác thực độc lập đảm bảo chỉ thiết bị hợp
    lệ mới được ghi dữ liệu và chỉ người dùng hợp lệ mới xem được.
  \item \textbf{Phòng Socket.io theo device/user:} Server emit đến room
    \texttt{device\_\{id\}} và \texttt{user\_\{id\}}, đảm bảo mỗi
    client chỉ nhận dữ liệu của thiết bị mình đang xem, không bị
    ``nhiễu'' từ các thiết bị khác.
\end{itemize}
